\appendixtitles{yes}
\appendixstart
\appendix
\section{Per-Cell Tables}\label{app:tables}

This appendix holds the per-cell tables that Sections~\ref{sec:results} and~\ref{sec:discussion}
cite. Each is generated from the results files, and every value in it is a macro of the generated
macro files or a count from the results files.

\begin{table}[H]
\caption{Every invalid window of a family cell, per arm, with its reasons. The count after a reason
is the number of windows that gave it; a window may give several.\label{tab:invalid}}
\centering\scriptsize
\begin{tabularx}{\textwidth}{@{}l l r X@{}}
\toprule
\textbf{Cell} & \textbf{Arm} & \textbf{Windows} & \textbf{Reasons} \\
\midrule
\tabinput{../results/tables/invalid}
\bottomrule
\end{tabularx}
\end{table}

\begin{table}[H]
\caption{The secondary metrics of the cost cells, one-port / dedicated. WL4 gives the server's CPU
per exchange and its resident memory at the end and at its peak. WL2 gives the time to first byte
at \PNinetyNine, in C3 only. Each is the median with its \CiPct\% interval. On W the CPU measure is
the server's processor cycles, read with \texttt{QueryProcessCycleTime}; a ratio of two arms needs no conversion.\label{tab:wl4}}
\begin{adjustwidth}{-\extralength}{0cm}
\centering\scriptsize
\begin{tabularx}{\fulllength}{@{}l l l X X X X@{}}
\toprule
\textbf{Hyp.} & \textbf{Host, backend} & \textbf{Protocol} & \textbf{CPU} & \textbf{Memory, end} &
\textbf{Memory, peak} & \textbf{TTFB, tail} \\
\midrule
\tabinput{../results/tables/wl4}
\bottomrule
\end{tabularx}
\end{adjustwidth}
\end{table}

\begin{table}[H]
\caption{The analysis clustered by lab job, which decides nothing: the cost cells that ran in more
than one job, all on W. Beside each clustered interval is the cell's interval by session.
\label{tab:byjob}}
\centering\scriptsize
\begin{tabularx}{\textwidth}{@{}l l l l X X@{}}
\toprule
\textbf{Hyp.} & \textbf{Host, backend} & \textbf{Protocol} & \textbf{Median} &
\textbf{Interval by job} & \textbf{Interval by session} \\
\midrule
\tabinput{../results/tables/by-job}
\bottomrule
\end{tabularx}
\end{table}

\begin{table}[H]
\caption{The secondary cells with sessions of their own, each with its valid sessions, median and
\CiPct\% interval, or whether it was not run or has no valid session. Section~\ref{sec:res-secondary}
gives the reasons.\label{tab:sec}}
\begin{adjustwidth}{-\extralength}{0cm}
\centering\scriptsize
\begin{tabularx}{\fulllength}{@{}l l X r l l@{}}
\toprule
\textbf{Cell} & \textbf{Host, backend} & \textbf{Ratio, metric} & \textbf{Valid} & \textbf{Median} &
\textbf{\CiPct\% interval} \\
\midrule
\tabinput{../results/tables/secondary}
\bottomrule
\end{tabularx}
\end{adjustwidth}
\end{table}

\begin{table}[H]
\caption{The secondary cells of the mechanisms and of B3. They give the time to first byte at a
fixed load for M1 and M3 and the CPU per connection of every mechanism cell. For M2 and M3 that
CPU is the front's. The last rows give B3 with the server in its other detection
mode.\label{tab:secmech}}
\begin{adjustwidth}{-\extralength}{0cm}
\centering\scriptsize
\begin{tabularx}{\fulllength}{@{}l l X r l l@{}}
\toprule
\textbf{Cell} & \textbf{Host, backend} & \textbf{Ratio, metric} & \textbf{Valid} & \textbf{Median} &
\textbf{\CiPct\% interval} \\
\midrule
\tabinput{../results/tables/secondary-mech}
\bottomrule
\end{tabularx}
\end{adjustwidth}
\end{table}

\begin{table}[H]
\caption{The parts of the competitor's footprint in each B3 cell, as medians over the valid sessions. $W$, $U$, $K_q$ and $K_s - \Kbase$ are in bytes per pending connection. ``Skb caches'' is the growth of the three
socket-buffer caches subtracted from $K_s$. ``Shared cache'' is the part of it in the merged cache,
whose co-tenants are \name{pool\_workqueue} and \SgPool. Both are in bytes over the host. $f$ is the
mean of the \texttt{f} field of \texttt{ss -tm} over the accepted sockets, in bytes, reported and not
added.\label{tab:b3parts}}
\begin{adjustwidth}{-\extralength}{0cm}
\centering\scriptsize
\begin{tabularx}{\fulllength}{@{}l l l X X X X X X X@{}}
\toprule
\textbf{Case} & \textbf{System} & \textbf{Backend} & $W$ & $U$ & $K_q$ & $K_s - \Kbase$ & \textbf{Skb caches} & \textbf{Shared cache} & $f$ \\
\midrule
\tabinput{../results/tables/b3-parts-competitor}
\bottomrule
\end{tabularx}
\end{adjustwidth}
\end{table}

\begin{table}[H]
\caption{The parts of the server's footprint in each B3 cell, in the units of Table~\ref{tab:b3parts}.\label{tab:b3partsserver}}
\begin{adjustwidth}{-\extralength}{0cm}
\centering\scriptsize
\begin{tabularx}{\fulllength}{@{}l l l X X X X X X X@{}}
\toprule
\textbf{Case} & \textbf{System} & \textbf{Backend} & $W$ & $U$ & $K_q$ & $K_s - \Kbase$ & \textbf{Skb caches} & \textbf{Shared cache} & $f$ \\
\midrule
\tabinput{../results/tables/b3-parts-server}
\bottomrule
\end{tabularx}
\end{adjustwidth}
\end{table}

\begin{table}[H]
\caption{Operations and copies of the cost cells, per connection in C1 and C3 and per request in C2,
from the valid windows (WL5). Each value is the mean over an arm's valid windows; ``Windows'' gives
them as one-port/dedicated. Operations are the system calls and ring submissions the counters name.
The bytes are payload bytes, copied in user space by the server's own code or received and sent
across the boundary. Both arms run replay and in-process dispatch. No byte was peeked in any
arm.\label{tab:ops-cost}}
\begin{adjustwidth}{-\extralength}{0cm}
\centering\scriptsize
\begin{tabularx}{\fulllength}{@{}l l l X r r r r r r r r@{}}
\toprule
& & & & \multicolumn{2}{c}{\textbf{Operations}} & \multicolumn{2}{c}{\textbf{Copied}} &
\multicolumn{2}{c}{\textbf{Received}} & \multicolumn{2}{c}{\textbf{Sent}} \\
\textbf{Hyp.} & \textbf{Host, backend} & \textbf{Protocol} & \textbf{Windows} & \textbf{One-port} &
\textbf{Ded.} & \textbf{One-port} & \textbf{Ded.} & \textbf{One-port} & \textbf{Ded.} &
\textbf{One-port} & \textbf{Ded.} \\
\midrule
\tabinput{../results/tables/ops-cost}
\bottomrule
\end{tabularx}
\end{adjustwidth}
\end{table}

\begin{table}[H]
\caption{Operations and copies per connection of M1's cells, one-port mode with in-process dispatch
in each detection mode, as means over the valid windows, in the units of
Table~\ref{tab:ops-cost}. ``Wakeups'' counts the wakeups during detection.\label{tab:ops-m1}}
\centering\scriptsize
\begin{tabularx}{\textwidth}{@{}l l l r r r r r r X@{}}
\toprule
\textbf{Host, backend} & \textbf{Protocol} & \textbf{Mode} & \textbf{Windows} & \textbf{Operations} &
\textbf{Copied} & \textbf{Received} & \textbf{Sent} & \textbf{Peeked} & \textbf{Wakeups} \\
\midrule
\tabinput{../results/tables/ops-m1}
\bottomrule
\end{tabularx}
\end{table}

\begin{table}[H]
\caption{Operations and copies per relayed connection of the server's relay on L, from \JobSt's
untimed windows, in front of the stub as in M3. The load is open-loop churn at \TraceRate{}
exchanges per second. Each value is the front's count over its process divided by the connections it relayed, in
the units of Table~\ref{tab:ops-cost}. The loop's waits and the accepts that find no connection
depend on the load.\label{tab:ops-relay}}
\centering\scriptsize
\begin{tabularx}{\textwidth}{@{}l l l r r r r r r X@{}}
\toprule
\textbf{Backend} & \textbf{Protocol} & \textbf{Mode} & \textbf{Relayed} & \textbf{Operations} &
\textbf{Copied} & \textbf{Received} & \textbf{Sent} & \textbf{Peeked} & \textbf{Wakeups} \\
\midrule
\tabinput{../results/tables/ops-relay}
\bottomrule
\end{tabularx}
\end{table}

\begin{table}[H]
\caption{The check of the server's system-call counters, one untimed window per cost cell of L in
each mode, from \JobSt. Each entry gives the row's checks that equal the count of the same calls,
by \texttt{perf stat} (the check) and by \texttt{perf trace -s} (its cross-check), over the row's
checks. The shortfall is \texttt{perf trace -s}'s largest, as a percentage of \texttt{perf stat}'s
count of a call, over the cell's two rows. C1 and C3 share one untimed load, open-loop churn, and
C2 runs keep-alive.\label{tab:counter-check}}
\centering\scriptsize
\begin{tabularx}{\textwidth}{@{}l l l r r r r X@{}}
\toprule
& & & \multicolumn{2}{c}{\textbf{Dedicated}} & \multicolumn{2}{c}{\textbf{One-port}} & \\
\textbf{Hyp.} & \textbf{Backend} & \textbf{Protocol} & \textbf{Stat} & \textbf{Trace} & \textbf{Stat} &
\textbf{Trace} & \textbf{Shortfall \%} \\
\midrule
\tabinput{../results/tables/counter-check}
\bottomrule
\end{tabularx}
\end{table}
